
\subsubsection{Motivating example}
\begin{wrapfigure}[3]{R}{6.5cm}
\centering
\vspace{-35pt}
\includegraphics[width=6.5cm]{"images/pointer-1.jpg"}
\caption{Region map motivating example}
\label{img-region-motivating-example}
\vspace{-15pt}
\end{wrapfigure}

\par

In the WP memory model named Typed, typed accesses with different type size
are considered to be in separated memory locations. Therefore, the WP Typed
memory model is unsound when applied to this example\footnote{If the architecture sets the
length of an \texttt{int} to be double the length of a \texttt{short}}:

\begin{lstlisting}[linewidth=6cm]
void g(short x, int y) {
    x = 1;
    short *p = (short *) &y;
    p[0] = 0;
    p[1] = 1;
}
\end{lstlisting}

One can read the \emph{region map} represented in Figure \ref{img-region-motivating-example} as:

\begin{itemize}
    \item White nodes are \emph{region labels} and represent variables.
    \item Region labels for arrays are formed as \texttt{variable[length]},
    for example \texttt{y[2]}.
    \item The upper pink region (whose label is \texttt{y[2]}) is accessed
    only by writing a \texttt{short} value. The representation \texttt{W(s)}
    is a compact form, meaning \texttt{W} for the write access and \texttt{s}
    specifies the C type \texttt{short}. The pink color is a redondant
    representation of a write access of a primitive type, as demonstrated for
    the region labelled \texttt{x}.
    \item The centered orange region, labelled \texttt{p}, is accessed as
    a pointer in read and write access. The representation \texttt{RW*} is
    a compact form meaning \texttt{R} for the read access, \texttt{W} for the
    write access, and \texttt{*} represents the pointer type.
    \item An arrow links the region \texttt{p} to the region \texttt{y[2]},
    meaning that the pointer \texttt{p} points to the region
    \texttt{y[2]}\footnote{More specifically, it points to some of its memory locations.}.
\end{itemize}
